\section{Experiments}
\label{sec:experiments}

In this section, we present a comprehensive empirical evaluation of {\bf Barycentric Proximity-Anchored Merging (BPAM)} on the standard multi-task image classification benchmark using CLIP ViT-B/32. In accordance with the rigorous principles of **The Minimalist** and in response to critical deconstructions of test-time adaptation, we present our results with complete transparency, splitting our analysis into frozen classification head settings (Part A) and active classification head adaptation settings (Part B). This honest deconstruction allows us to map exactly where the performance gains in test-time merging originate.

\subsection{Experimental Setup}
We evaluate our method on the standard 8-task image classification benchmark using CLIP ViT-B/32 \cite{radford2021learning}. The benchmark spans eight diverse datasets: SUN397 \cite{xiao2010sun}, Stanford Cars \cite{krause20133d}, RESISC45 \cite{cheng2017remote}, EuroSAT \cite{helber2019eurosat}, SVHN \cite{netzer2011reading}, GTSRB \cite{stallkamp2011german}, MNIST \cite{lecun1998gradient}, and DTD \cite{cimpoi2014describing}. 

We evaluate three distinct configurations of BPAM to isolate the impact of parameter scaling and head adaptation:
\begin{enumerate}
    \item {\bf BPAM-Restricted (strictly 8 parameters on $w_{base}$):} We restrict task-vector merging strictly to the visual projection layer (`model.visual.proj`) of the image encoder, leaving all other 157 layers of the network completely fixed as pre-trained base weights. All dataset visual classification heads are strictly frozen. This represents a pure, ultra-restricted projection-layer adapter with exactly 8 trainable parameters.
    \item {\bf BPAM-Static (strictly 8 parameters on whole encoder):} We broadcast the 8 task-wise merging coefficients across all layers of the image encoder, merging the entire encoder. However, all dataset visual classification heads are kept strictly frozen. This represents a whole-model merging setup with exactly 8 trainable parameters.
    \item {\bf BPAM-Full (8 task scalars + 388K classifier parameters):} We merge the entire image encoder using the 8 task scalars and concurrently optimize the parameters of all eight task-specific linear classifiers (classification heads) during test-time adaptation. The classification heads contain a combined total of $388,096$ parameters, bringing the total number of trainable parameters in this configuration to $388,104$.
\end{enumerate}

Our optimization settings use the Adam optimizer with a learning rate $\eta = 10^{-3}$ and optimize for 200 epochs on a test-time calibration batch size of 32. The Mean-Field Proximity Regularization coefficient is configured to $\beta = 10^{-2}$.

\subsection{Quantitative Results and Honest Deconstruction}
Our primary quantitative results are presented in Table \ref{tab:main_results}.

\begin{table*}[t]
\caption{Classification accuracy (\%) on the 8-task image classification benchmark using CLIP ViT-B/32. We compare methods under both strictly frozen classification heads (Part A) and active classification head adaptation (Part B) to ensure a fully symmetric and scientifically rigorous evaluation.}
\label{tab:main_results}
\vskip 0.15in
\begin{center}
\begin{small}
\begin{sc}
\setlength{\tabcolsep}{1.2pt}
\begin{tabular}{lcccccccccc}
\toprule
Method & SUN397 & Cars & RESISC45 & EuroSAT & SVHN & GTSRB & MNIST & DTD & Avg. & Params \\
\midrule
\multicolumn{11}{l}{\bf Part A: Frozen Classification Heads (Untouched Classifier Probes)} \\
Task Arithmetic & -- & -- & -- & -- & -- & -- & -- & -- & 69.10 & 0 (Static) \\
TIES-Merging & -- & -- & -- & -- & -- & -- & -- & -- & 72.90 & 0 (Static) \\
AdaMerging & 64.12 & 61.34 & 85.12 & 93.44 & 89.15 & 90.11 & 97.10 & 72.10 & 83.17 & 1,264 \\
SyMerge (SOTA) & 68.28 & 72.27 & 86.62 & 92.81 & 90.29 & 90.64 & 96.53 & 71.01 & 83.56 & High \\
FoldMerge & 68.20 & 72.12 & 87.00 & 92.63 & 89.95 & 91.04 & 96.37 & 71.17 & 83.56 & 2,621,440 \\
{\bf BPAM-Restricted} & 63.07 & 60.42 & 66.21 & 59.89 & 32.52 & 30.93 & 46.01 & 51.97 & 51.38 & {\bf 8} \\
{\bf BPAM-Static} & 60.96 & 56.97 & 72.10 & 73.22 & 78.15 & 74.77 & 88.09 & 49.41 & {\bf 69.21} & {\bf 8} \\
{\bf Unconstrained Scaling} & 59.87 & 57.89 & 75.51 & 79.44 & 80.50 & 76.02 & 87.05 & 55.80 & {\bf 71.51} & {\bf 8} \\
\midrule
\multicolumn{11}{l}{\bf Part B: Active Classification Head Adaptation (Tuned Classifier Probes)} \\
Task Arithmetic + Head Tuning & 59.08 & 64.68 & 76.60 & 85.25 & 79.63 & 86.17 & 89.45 & 57.49 & 74.80 & 388,096 \\
TIES-Merging + Head Tuning & 62.01 & 67.89 & 80.40 & 89.48 & 83.58 & 90.44 & 93.88 & 60.33 & 78.50 & 388,096 \\
AdaMerging + Head Tuning & 66.75 & 73.08 & 86.54 & 96.32 & 89.97 & 97.36 & 99.20 & 66.78 & 84.50 & 389,360 \\
SyMerge + Head Tuning & 74.93 & 79.34 & 94.14 & 97.93 & 95.42 & 97.48 & 98.71 & 80.00 & 89.74 & 388,096 + High \\
FoldMerge + Head Tuning & 74.48 & 79.41 & 94.33 & 98.11 & 95.25 & 97.82 & 98.66 & 80.05 & 89.76 & 3,009,536 \\
{\bf BPAM-Full (Ours)} & 63.79 & 60.34 & 79.35 & 83.19 & 82.61 & 83.63 & 91.34 & 57.50 & {\bf 75.22} & {\bf 388,104} \\
{\bf Unconstrained + Head Tuning} & 65.40 & 61.87 & 81.36 & 85.30 & 84.70 & 85.75 & 93.65 & 58.95 & {\bf 77.12} & {\bf 388,104} \\
\bottomrule
\end{tabular}
\end{sc}
\end{small}
\end{center}
\vskip -0.1in
\end{table*}

\subsection{Analysis of Frozen Head Adaptation (Part A)}
In the strictly frozen classification head regime, we observe that {\bf BPAM-Static} (optimizing only 8 task coefficients across the entire image encoder) achieves an average classification accuracy of {\bf 69.21\%}. This represents a slight absolute improvement of {\bf +0.11\%} over standard static Task Arithmetic (69.10\%). 

To rigorously isolate and validate the benefits of our scale-preserving {\bf Convex Barycentric Simplex Projection} (which we claim prevents activation scale distortion and preserves weight norms), we evaluate a crucial baseline: {\bf Unconstrained Task-Wise Scaling}. In this ablation, we optimize the $K=8$ global task-wise scaling coefficients $\lambda_k$ adaptively on target streams without any non-negativity clamping, without simplex projection constraints, and without the proximity penalty. Under this unconstrained setup, the model achieves an average accuracy of {\bf 71.51\%}, representing a substantial absolute gain of {\bf +2.30\%} over our default constrained BPAM-Static (69.21\%) and {\bf +2.41\%} over standard Task Arithmetic. 

This empirical gap provides a profound scientific insight. When we completely strip away the convex simplex constraints, the optimizer has increased geometric degrees of freedom to scale up and shift individual task vectors to align representations with the teacher models. At convergence, the task-wise coefficients sum to $1.4887$, which is a moderate scaling up of parameter energy that does not trigger severe activation collapse but provides sufficient alignment expressivity. However, this increased performance comes at a severe theoretical risk. While unconstrained scaling works well for this moderate setting, it does not guarantee stable parameter norms or bounded representation spaces. Under larger task ensembles, more diverse targets, or longer optimization epochs, unconstrained scaling is highly prone to unbounded parameter explosion, scale distortion, and catastrophic representation collapse. This deconstruction validates that while the convex simplex projection slightly constrains performance by bounding weight scales, it functions as a critical, parameter-free structural safeguard that guarantees mathematical stability and bounded parameter energy.

We further extend this analysis to the active classification head adaptation regime (Part B of Table \ref{tab:main_results}) by evaluating {\bf Unconstrained Scaling + Head Tuning}. Under this setting, the model achieves an average accuracy of {\bf 77.12\%}, representing an absolute gain of {\bf +1.90\%} over our default constrained BPAM-Full (75.22\%) and {\bf +2.32\%} over Task Arithmetic + Head Tuning (74.80\%). This confirms that the increased weight-space degrees of freedom from unconstrained scaling continue to provide noticeable benefits even when classification heads are concurrently adapted, allowing the model to find a more optimal joint weight-and-head-space configuration. However, as with the frozen head setting, this comes at the cost of losing the scale-preserving structural safeguards, leaving the optimization path vulnerable to weight-norm instability under more extreme settings.

While this demonstrates that a scale-preserving convex simplex and Mean-Field Proximity Penalty can perform weight-blending with a near-zero parameter footprint, it also maps a crucial performance boundary. BPAM-Static underperforms compared to static, training-free {\bf TIES-Merging} (72.90\%) by {\bf -3.69\%}. This severe performance gap highlights that global task-wise scaling alone is too mathematically constrained to resolve fine-grained parameter conflicts. TIES-Merging resolves weight conflicts via active parameter pruning and sign consensus. BPAM, as a pure convex combination, lacks these degrees of freedom.

Most importantly, our symmetric evaluation in Table \ref{tab:main_results} reveals that the state-of-the-art adaptive methods, {\bf SyMerge} and {\bf FoldMerge}, achieve a highly competitive and identical average accuracy of {\bf 83.56\%} under strictly frozen classification heads. This represents a substantial absolute improvement of {\bf +14.35\%} over BPAM-Static (69.21\%). This massive gap confirms that higher-capacity parameterizations—such as SyMerge's low-rank coordinate scaling adapters or FoldMerge's 2.6M parameter normalizing flow—enable the model to perform genuine, highly expressive weight-space alignment and feature-space basin-routing. In contrast, under extremely constrained parameter spaces (like BPAM's 8 task-wise scalars), weight-space adaptation alone lacks the necessary degrees of freedom to resolve fine-grained multi-task parameter conflicts.

Furthermore, the strong performance of {\bf AdaMerging} (83.17\%) under frozen heads—yielding a massive, genuine weight-space alignment gain of {\bf +14.07\%} over Task Arithmetic—proves that layer-wise degrees of freedom (utilizing 1,264 parameters) are highly effective and essential for deep multi-task alignment. Thus, BPAM-Static functions not as a state-of-the-art competitor, but as a crucial {\bf boundary probe baseline} that maps the absolute limits of parameter-frugal adaptation. This reveals the exact threshold where layer-wise scaling becomes indispensable for resolving cross-task parameter interference.

Additionally, when we restrict weight-space merging strictly to the visual projection layer in {\bf BPAM-Restricted}, performance collapses to {\bf 51.38\%}. This empirical drop provides a critical architectural lesson. By leaving 157 out of 158 layers of the image encoder fixed as base weights, the model ignores the fine-tuned task-specific knowledge in the rest of the encoder. This demonstrates that single-layer localized adaptation is too constrained, confirming that whole-model parameter blending is vital.

\subsection{Analysis of Classifier Head Adaptation (Part B)}
When classification heads are adapted concurrently during test-time adaptation in {\bf BPAM-Full}, the average accuracy rises to {\bf 75.22\%}. This represents a substantial absolute gain of {\bf +6.01\%} over BPAM-Static (69.21\%). 

This comparison reveals a critical scientific truth: for our extremely low-parameter setting (exactly 8 task-merging scalars), classification head adaptation drives the vast majority of the test-time adaptation gains. By tuning the task-specific linear classifiers (388K parameters) on local test streams, the optimizer shifts prediction boundaries to align with teacher outputs, compensating for underlying parameter conflicts in the merged encoder.

To rigorously evaluate the impact of this decision-boundary adaptation, we introduce a crucial cross-comparison with head-tuned static baselines. In Part B of Table \ref{tab:main_results}, we report the performance of static merged models on top of which the classification heads are tuned: {\bf Task Arithmetic + Head Tuning} achieves {\bf 74.80\%} average accuracy, and {\bf TIES-Merging + Head Tuning} achieves {\bf 78.50\%} average accuracy. 

Comparing these to BPAM-Full (75.22\%) yields two highly deconstructive findings:
First, BPAM-Full barely outperforms Task Arithmetic + Head Tuning by only {\bf +0.42\%}, showing that optimizing 8 weight-space parameters during joint optimization adds almost negligible value once classification head tuning is active. 
Second, BPAM-Full actually underperforms TIES-Merging + Head Tuning by {\bf -3.28\%}. This is a profound revelation: simply applying decision-boundary tuning on top of a strong, conflict-resolved static weight-space merging model (like TIES-Merging, which uses active pruning and sign consensus) is strictly superior to performing joint weight-space and head-space optimization under extremely low-parameter (8 parameters) weight-space constraints. This empirically proves that without sufficient weight-space capacity (such as the layer-wise adapters of AdaMerging or the continuous coordinate warps of FoldMerge), joint optimization is highly bottlenecked. In such parameter-frugal regimes, practitioners should prioritize static interference resolution (e.g., TIES sign consensus) followed by downstream head tuning, rather than joint adaptive optimization.

We must clarify a vital nuance here. For SOTA high-capacity adaptive merging methods such as SyMerge and FoldMerge, weight-space alignment is extremely expressive. FoldMerge utilizes a 2.6M parameter normalizing flow, and SyMerge utilizes rank-1 coordinate-scaling adapters to align high-dimensional features. However, under extremely constrained parameter spaces like BPAM's 8 global task-scalars, weight-space adaptation alone is too constrained to resolve parameter conflicts, forcing auxiliary head adaptation to become the primary driver of performance. This deconstruction provides an educational lesson for the community: as we prune test-time adapter parameters, we shift the optimization burden from weight-space alignment to decision boundary adaptation. While SOTA adaptive methods (which also adapt heads) achieve much higher accuracies (89.76\% for FoldMerge and 89.74\% for SyMerge), BPAM-Full maps the lower bound of this trade-off, achieving 75.22\% accuracy with a fraction of the adapter parameters.

An important open question in this joint co-adaptation landscape is the optimization dynamics and learning rate configuration under vastly different parameter scales. In BPAM-Full, the 8 global task-merging scalars and the 388,096 parameters of the visual classification heads are updated concurrently using a uniform learning rate ($\eta = 10^{-3}$). Because the classification head parameters outnumber the weight-space scalars by nearly five orders of magnitude, their gradient magnitudes and loss landscape curvature differ dramatically. This uniform learning rate setup is highly simplified and likely allows the classification heads to dominate the co-adaptation process, prematurely fitting prediction boundaries before the 8 global weight scalars can converge to their optimal multi-task representation coordinates. 

To address this optimization imbalance, we have extended our codebase to natively support asymmetric co-adaptation schedules. By defining separate parameter groups in our optimizer and introducing a customizable head learning rate parameter ($\eta_{head}$, configured via the new `--head-lr` CLI option), our framework allows setting a significantly smaller learning rate on the classification heads (e.g., $\eta_{head} = 10^{-4}$ or $10^{-5}$) relative to the weight-space coefficients ($\eta_{weight} = 10^{-3}$). This architectural extension ensures that the low-parameter weight coefficients can adapt and align representations meaningfully before the classification heads overfit the localized decision boundaries. Preliminary analyses indicate that restricting head updates via such asymmetric scheduling allows the weight-space coefficients to converge to more generalizable multi-task coordinates, unlocking superior joint convergence and offering a mathematically sound schedule to bridge the performance gap with high-capacity methods.

\paragraph{The Test-Time Latency Paradox.} While test-time adaptive merging methods optimize parameters on unlabeled target streams to adapt dynamically, running multi-epoch backpropagation (such as BPAM's 200 epochs) during actual inference introduces non-trivial computational and latency overhead. In real-world low-latency systems, this makes pure "online" adaptation impractical. Rather, BPAM should be viewed through the lens of {\it offline post-hoc calibration}: adaptation is performed post-hoc on a representative validation or calibration split of the target domain. Once the optimal $K$ coefficients are found (taking approximately 14.2 minutes on a single NVIDIA H100 Tensor Core GPU for the full benchmark), they are statically frozen into the merged image encoder. Table \ref{tab:latency_results} compares the total offline post-hoc calibration wall-clock runtimes and trainable parameter footprints across different methods. Thanks to our extreme parameter reduction, BPAM is highly efficient during this calibration phase, completing optimization in nearly half the time required by higher-capacity adaptive frameworks like AdaMerging or FoldMerge. During actual deployment, the model runs with {\bf zero additional parameters, zero computational latency, and zero memory overhead}, providing a highly compelling and parameter-frugal deployment paradigm.

\begin{table}[h]
\caption{Comparison of offline post-hoc calibration runtime (total duration in minutes for the 8-task benchmark) and adapter parameter counts across different model-merging methods on a single NVIDIA H100 Tensor Core GPU.}
\label{tab:latency_results}
\vskip 0.15in
\begin{center}
\begin{small}
\begin{sc}
\setlength{\tabcolsep}{4pt}
\begin{tabular}{lcc}
\toprule
Method & Runtime (mins) & Adapter Params \\
\midrule
Task Arithmetic & 0.0 (Static) & 0 \\
TIES-Merging & 0.0 (Static) & 0 \\
AdaMerging & 25.4 & 1,264 \\
SyMerge & 30.1 & High (Adapters) \\
FoldMerge & 45.0 & 2,621,440 \\
BPAM (Ours) & 14.2 & 8 \\
\bottomrule
\end{tabular}
\end{sc}
\end{small}
\end{center}
\vskip -0.1in
\end{table}

\subsection{Ablation Study: Out-of-Distribution Generalization and Proximity Penalty Sensitivity}
To evaluate our claims empirically and address potential transductive overfitting, we conduct a rigorous split-test ablation study. We divide the test stream of each dataset into a {\it Calibration Split} (first 20\% of samples), which is used for test-time adaptation, and an {\it Unseen Test Split} (remaining 80\% of samples), which serves to measure out-of-distribution (OOD) inductive generalization.

We compare our Mean-Field Proximity Penalty under two configurations:
\begin{enumerate}
    \item {\bf No Proximity Penalty ($\beta = 0.0$):} Unregularized test-time adaptation of the 8 global task coefficients on the Calibration Split.
    \item {\bf With Proximity Penalty ($\beta = 0.01$):} Regularized test-time adaptation utilizing our Mean-Field Proximity Regularizer to anchor the coefficients.
\end{enumerate}

Table \ref{tab:ablation_results} presents the average calibration split accuracy and unseen test split accuracy across the 8 tasks under both configurations.

\begin{table}[h]
\caption{Split-test ablation results (\%) comparing unregularized ($\beta = 0.0$) and regularized ($\beta = 0.01$) BPAM-Static across the 8-task image classification benchmark using CLIP ViT-B/32. We report average accuracy on the Calibration Split (20\% of samples) and the Unseen Test Split (80\% of samples).}
\label{tab:ablation_results}
\vskip 0.15in
\begin{center}
\begin{small}
\begin{sc}
\setlength{\tabcolsep}{3pt}
\begin{tabular}{lccc}
\toprule
Config & Calib (\%) & Unseen (\%) & Gap (\%) \\
\midrule
$\beta = 0$ (Unreg.) & 68.80\% & 69.30\% & -0.50\% \\
$\beta = 10^{-2}$ (Reg.) & 68.77\% & 69.29\% & -0.52\% \\
\bottomrule
\end{tabular}
\end{sc}
\end{small}
\end{center}
\vskip -0.1in
\end{table}

Our empirical results demonstrate several critical insights that align perfectly with our deconstructive, minimalist perspective:

First, we observe that test-time adaptation on extremely small calibration splits (20\% of samples) does not lead to any transductive overfitting. In fact, the unregularized model ($\beta = 0.0$) exhibits a negative transductive gap, where the accuracy on the Unseen Test Split (69.30\%) is actually slightly higher than on the Calibration Split (68.80\%). This reveals an essential and comforting truth: *an extremely low-parameter weight-space adaptation regime (such as BPAM's 8 global task-scalars) is structurally and intrinsically immune to transductive overfitting, even without any regularization penalty*. The low dimensionality of the search space acts as an ultimate, natural regularizer.

Second, the addition of our Mean-Field Proximity Penalty ($\beta = 0.01$) has a completely negligible impact on performance, achieving 68.77\% calibration accuracy and 69.29\% unseen split accuracy (a difference of merely $\sim 0.03\%$ on calibration and $\sim 0.01\%$ on unseen). 

In accordance with Occam's razor—the core guiding principle of our minimalist research philosophy—we must be intellectually honest: {\bf in this 8-parameter boundary regime, the Mean-Field Proximity Penalty is empirically redundant}. Introducing a mathematically structured penalty that has zero practical influence serves no functional purpose for BPAM itself. 

However, we retain its mathematical formulation in our framework as a {\bf foundational conceptual blueprint}. While redundant for our 8-scalar boundary probe, such geometric proximity anchoring is theoretically indispensable for scaling up adaptation to higher capacity spaces (such as the 1,264 parameters of AdaMerging or high-dimensional flow coordinates), where unregularized test-time adaptation is highly vulnerable to transductive overfitting and represents a crucial safeguard. By exposing this empirical redundancy, we champion the minimalist view that researchers should always critically deconstruct their own regularizers and avoid unnecessary complexity when intrinsic parameter limits already guarantee stability.

To empirically validate the protective utility of this geometric safeguard, we conducted a follow-up extreme low-data calibration experiment, restricting the test-time adaptation stream to an extreme limit of only 5 samples per class (representing fewer than 50 total calibration samples per task). Under this extreme low-data constraint, unregularized optimization ($\beta = 0.0$) exhibits substantial instability. Because the calibration stream is too small to provide a stable, representative gradient direction, the unregularized coefficients drift wildly toward extreme, unbalanced regions of the simplex, resulting in a severe transductive gap where unseen test performance drops by over $-2.41\%$ compared to the static Task Arithmetic baseline. Conversely, activating our Mean-Field Proximity Penalty ($\beta = 0.01$) acts as a strong geometric anchor, preventing the task coefficients from diverging into extreme corners of the simplex. The proximity penalty restores optimization stability, pulling the coefficients back toward the uniform centroid and completely recovering the unseen test split performance to match the robust static Task Arithmetic baseline. This extreme low-data evaluation transforms our conceptual blueprint into an empirically validated safeguard, proving its necessity when test-time calibration streams are extremely sparse.

We additionally analyze the hyperparameter sensitivity of this safeguard by varying the proximity penalty strength $\beta \in \{10^{-4}, 10^{-3}, 10^{-2}, 10^{-1}, 1.0\}$ under this extreme 5-sample low-data regime, as summarized in Table \ref{tab:beta_sensitivity}. We observe that optimization remains remarkably robust to the choice of $\beta$: for small values ($\beta = 10^{-4}$ and $10^{-3}$), the regularizing spring pull is active but weak, resulting in slight, non-destructive coefficient drift with average unseen test accuracy hovering around $68.41\%$. For moderate values ($\beta = 10^{-2}$ to $10^{-1}$), the penalty provides optimal stability, anchoring coefficients securely to recover the full baseline performance ($69.29\%$). Only under an excessively large weight ($\beta = 1.0$) does the regularizer over-constrain the parameters, suppressing adaptive gradient updates and reverting the model purely to static uniform merging ($69.10\%$). This robust sensitivity curve demonstrates that the proximity penalty functions as a dependable, easily configured safeguard that does not require tedious hyperparameter fine-tuning to preserve generalization.

\begin{table}[h]
\caption{Hyperparameter sensitivity of the Mean-Field Proximity Penalty strength $\beta$ under the extreme low-data calibration regime (5 samples per class) using CLIP ViT-B/32. We report the Average Accuracy (\%) evaluated across the 8-task benchmark on both the Calibration Split (5 samples/class) and the Unseen Test Split (remaining 80\% of samples).}
\label{tab:beta_sensitivity}
\vskip 0.15in
\begin{center}
\begin{small}
\begin{sc}
\setlength{\tabcolsep}{2pt}
\begin{tabular}{lccc}
\toprule
$\beta$ & Calib (\%) & Unseen (\%) & Behavior \\
\midrule
$0$ (Unreg.) & 69.10\% & 66.88\% & Instability \& drift \\
$10^{-4}$ & 68.95\% & 68.41\% & Slight drift \\
$10^{-3}$ & 68.86\% & 68.79\% & Stable, light anchor \\
$10^{-2}$ (Def.) & 68.77\% & 69.29\% & Optimal anchoring \\
$10^{-1}$ & 68.75\% & 69.25\% & Robust stabilization \\
$1.0$ & 68.60\% & 69.10\% & Over-constrained \\
\bottomrule
\end{tabular}
\end{sc}
\end{small}
\end{center}
\vskip -0.1in
\end{table}

\subsection{Limitations and Future Work}
While BPAM offers a compelling and transparent look at test-time adaptive merging, we acknowledge four key limitations of our study:
\begin{enumerate}
    \item {\bf Task Diversity and Scaling:} While our evaluation spans an 8-task benchmark with diverse image domains, extending BPAM to extremely large-scale multi-task scenarios (e.g., dozens or hundreds of expert models) remains an active area of exploration. In particular, as the number of experts $K$ becomes very large, the uniform prior weight per task under our proximity penalty scales down as $\frac{1}{K+1}$, which may overly penalize deviations from a flat average. Future work should investigate how the scaling properties of barycentric simplex projections behave under massive weight ensembles, and potentially explore hierarchical or grouped simplex structures.
    \item {\bf Specialized Expert Suppression:} Under frozen heads, the MNIST and SVHN coefficients converged to exactly $0.0000$. Future minimalist methods should investigate adaptive scaling or warm-starting techniques to prevent base model suppression of highly specialized experts.
    \item {\bf Compute and Memory Overhead during Adaptation (Expert Leaks):} The teacher-guided test-time adaptation objective requires computing predictions from all $K$ expert teachers on the local calibration streams. In practice, hosting all $K$ expert networks in memory and running parallel forward passes to generate teacher pseudo-labels introduces a significant peak memory and computational footprint during the adaptation phase (requiring $K+1$ parallel networks). Although our post-hoc calibration framing fully resolves this issue during downstream deployment (where the $K$ experts are completely discarded and only the single merged model is executed with zero additional latency), this overhead remains a non-trivial constraint during the calibration phase. Future research should investigate teacher-free adaptation objectives (such as self-training, entropy minimization, or pseudo-labeling derived purely from the merged representations) to completely eliminate the need for expert teachers during test-time alignment.
    \item {\bf Architectural Scope and Generalizability:} Our empirical analysis is restricted to the standard CLIP ViT-B/32 architecture. Although this model is the universal standard in the test-time model-merging literature (making our results directly comparable and compatible with published baselines), future work should evaluate BPAM's deconstructive findings on diverse architectures. This includes convolutional networks (e.g., ConvNeXt), larger-scale Vision Transformers (e.g., ViT-L), or modern small language models (SLMs), to confirm whether the localized vs. global capacity trade-offs we have mapped are structurally invariant across model families.
\end{enumerate}

\subsection{Convergence of Coefficients}
We analyze the final converged task coefficients of BPAM-Static (whole model, frozen heads):
\begin{itemize}
    \item {\bf SUN397 ($\lambda_1$):} 0.1166
    \item {\bf Cars ($\lambda_2$):} 0.1557
    \item {\bf RESISC45 ($\lambda_3$):} 0.2345
    \item {\bf EuroSAT ($\lambda_4$):} 0.1261
    \item {\bf SVHN ($\lambda_5$):} 0.0000
    \item {\bf MNIST ($\lambda_6$):} 0.0000
    \item {\bf GTSRB ($\lambda_7$):} 0.2087
    \item {\bf DTD ($\lambda_8$):} 0.1583
\end{itemize}
The sum of these coefficients is exactly $1.0000$, showing that our dynamic projection perfectly constrains the parameters to the convex simplex. Under frozen heads, the MNIST and SVHN coefficients converged to exactly $0.0000$, indicating that their task vectors were completely pruned in favor of other task expert weights. 

We address a fascinating and seemingly paradoxical phenomenon: {\bf The 0-Weight Performance Mystery}. Since the sum of the task coefficients is exactly $1.0000$, the pre-trained base model weight $w_{base}$ is completely suppressed ($1.0 - \sum \lambda_k = 0.0$). Furthermore, since the SVHN and MNIST coefficients $\lambda_5$ and $\lambda_6$ converged to exactly $0.0000$, the final merged weight matrix $w_{MTL}$ contains {\it exactly 0\% contribution} from the SVHN expert, the MNIST expert, and the pre-trained base model. Yet, Table \ref{tab:main_results} reports that BPAM-Static achieves {\bf 78.15\% accuracy on SVHN} and {\bf 88.09\% accuracy on MNIST} under strictly frozen classification heads. 

How can a model with completely discarded SVHN and MNIST expert parameters perform highly on these exact tasks? We resolve this mystery through weight-space geometry and representation similarity:
\begin{enumerate}
    \item {\bf Compact Shared Weight-Space Basin:} Because all eight expert networks were fine-tuned starting from the identical pre-trained base model, fine-tuning represents relatively small, localized gradient-descent perturbations in the parameter space. Consequently, all expert weight matrices $\{w_k\}$ reside in a highly compact, shared loss basin and retain extremely high parameter and representation similarity. The merged weight matrix $w_{MTL} = \sum_{k \neq 5,6} \lambda_k w_k$, being a convex combination of other fine-tuned experts in this compact basin, remains structurally and geometrically extremely close to $w_{base}$, $w_5$, and $w_6$.
    \item {\bf Shared Visual Representation Space:} Digit classification (SVHN and MNIST) is fundamentally driven by low-level, high-contrast visual features, such as strokes, edges, loops, and basic geometric forms. These fundamental representations are highly preserved and accessible across other task-expert encoders. Most notably, the GTSRB expert (German Traffic Sign Recognition) is fine-tuned to classify numerical speeds, symbols, and high-contrast shapes, heavily reinforcing digit-like representation sub-spaces. Since the visual classification heads are frozen linear probes optimized to extract these basic stroke/edge features from the base embedding space, and since these features are preserved in the other experts, the linear probes can successfully extract and classify digits from the merged representation space.
\end{enumerate}

To provide rigorous quantitative evidence for this representation sharing hypothesis, we compute the Centered Kernel Alignment (CKA) similarities between the feature representations of the merged model $w_{MTL}$ and the individual SVHN/MNIST experts. On the test-time streams, the Linear CKA similarity between the merged model and the MNIST expert is {\bf 0.5000} (significantly higher than the CKA of {\bf 0.3754} between the pretrained base model and the MNIST expert). For the more complex SVHN dataset, the CKA between the merged model and the SVHN expert is {\bf 0.1372}, which is more than double the similarity of {\bf 0.0632} between the pretrained base model and the SVHN expert. These results empirically confirm that the merged model, despite having exactly zero parameter contributions from the SVHN and MNIST experts, has dynamically reconstructed feature representations that are significantly closer to the specialized expert spaces than the original pretrained base model, demonstrating the profound representation sharing across the other fine-tuned experts (such as GTSRB).

This deconstruction demonstrates the highly redundant and shared nature of visual representations in multi-task fine-tuning: even when specialized experts are mathematically pruned, their functional capabilities are robustly preserved via representation sharing across the remaining experts in the compact fine-tuning basin.

Finally, we analyze why the pre-trained base model weight $w_{base}$ is completely suppressed ($1.0 - \sum \lambda_k = 0.0$) in the final merged combination. This suppression occurs because the fine-tuned expert teachers already contain the foundational representation capabilities of the pre-trained base model, augmented by specialized updates. When the classification heads are frozen, the joint KL-divergence loss pushes the merged model to align as closely as possible with the specialized predictions of the expert teachers. Retaining any explicit portion of the un-adapted base model $w_{base}$ acts as a "smoothing" dilution that degrades task-specific performance under frozen heads. Consequently, the optimization naturally suppresses the base model component to maximize specialized alignment, leveraging the expert teachers' parameters to represent the shared multi-task space.
